\pdfinfoomitdate=1
\pdftrailerid{}
\pdfsuppressptexinfo=15
\newcommand{\DoNotLoadEpstopdf}{}
\documentclass[11pt]{article}
\usepackage[T1]{fontenc}
\usepackage[utf8]{inputenc}
\usepackage{lmodern,amsmath,amssymb,amsthm,mathtools,booktabs,array}
\usepackage[margin=1in]{geometry}
\usepackage{microtype}
\usepackage[hidelinks]{hyperref}
\usepackage{enumitem}
\setlist{itemsep=3pt,topsep=5pt}
\allowdisplaybreaks[2]
\emergencystretch=2em
\newtheorem{theorem}{Theorem}[section]
\newtheorem{lemma}[theorem]{Lemma}
\newtheorem{proposition}[theorem]{Proposition}
\newtheorem{corollary}[theorem]{Corollary}
\theoremstyle{remark}\newtheorem{remark}[theorem]{Remark}
\newcommand{\eps}{\varepsilon}
\newcommand{\OO}{\mathcal O}
\newcommand{\R}{\mathbb R}
\newcommand{\C}{\mathbb C}
\newcommand{\ii}{\mathrm i}
\newcommand{\dd}{\,\mathrm d}
\newcommand{\G}{\mathcal G}
\newcommand{\B}{\mathcal B}
\DeclareMathOperator{\Log}{Log}
\hypersetup{pdftitle={Report187 Hidden Oscillations in Square-Root Factorial Sampling},pdfauthor={Research report},pdfsubject={OEIS A326805 exact harmonics all fixed orders sharp oscillations and inverse thresholds}}
\title{Report187\\[5pt]Hidden Oscillations in\\Square-Root Factorial Sampling}
\author{OEIS A326805\\Exact harmonics, all fixed orders, and inverse thresholds}
\date{3 October 2026}
\begin{document}
\maketitle
\begin{abstract}
For the square-root sampler
\[
 S(x)=\sum_{k\ge0}\frac{x^{\sqrt k}}{\Gamma(1+\sqrt k)},\qquad a(n)=\lfloor S(n)\rfloor,
\]
we prove the OEIS conjecture $a(n)\sim2ne^n$ and resolve the much smaller, but unbounded, discrepancy. Its sharp envelope is $n^{1/8}(\log n)^{3/8}$, with an explicit Lambert--$W$ chirp phase and a computable expansion to every fixed logarithmic order. An exact finite-harmonic formula controls all omitted modes together, independently of the fixed-mode asymptotics. We prove the complete interval of subsequential limits, a complex-sector bound, a convergent Lagrange expansion using the exact discrepancy, and exact integer threshold identities. The naive threshold $\lceil W(Y/2)\rceil$ fails infinitely often in both directions, despite an exponentially small real inverse displacement. The accompanying offline source archive provides exact rational coefficient algebra, rational truncation-tail checks, corruption tests, separate optional numerical diagnostics, and a deterministic PDF builder.
\end{abstract}

\paragraph{Scope and attribution.}
The leading equivalent was conjectured by V\'aclav Kot\v e\v sovec in the OEIS record in 2019. Abel--Plana summation, complex Stirling expansions, the continuous Volterra--Ramanujan identity and analytic Lagrange inversion are classical inputs. The subject developed here is their target-specific application to the hidden sampling discrepancy, its harmonic hierarchy, and its inverse consequences. The expansion is additive on a positive envelope, never relative to a sine near its zeros. No uniform growing-harmonic expansion, convergent infinite asymptotic harmonic sum, effective asymptotic onset, interval certification of floating evaluations, or worldwide priority claim is made.

\newpage
{\small\tableofcontents}
\newpage

\section{The model and the main statements}\label{sec:main}
For $x>0$ set
\begin{equation}\label{eq:model}
 S(x)=\sum_{k=0}^{\infty}\frac{x^{\sqrt k}}{\Gamma(1+\sqrt k)},\qquad
 M(x)=2xe^x,\qquad D(x)=S(x)-M(x).
\end{equation}
At zero set $S(0)=1$ by continuity. The integer sequence $a(n)=\lfloor S(n)\rfloor$, $n\ge0$, is OEIS A326805 \cite{oeis}. Its first ten values are
\[
1,\ 6,\ 29,\ 120,\ 436,\ 1484,\ 4841,\ 15352,\ 47695,\ 145855.
\]
The continuous integral suggested by the sampling density $\dd k=2u\dd u$ is
\begin{equation}\label{eq:I}
 I(x)=2\int_0^\infty \frac{u x^u}{\Gamma(1+u)}\dd u.
\end{equation}
It differs from $M(x)$ by only $\OO((\log x)^{-2})$. Discrete sampling creates a substantially larger oscillatory error.

Throughout, $W$ denotes the positive real Lambert function: $W(t)e^{W(t)}=t$ for $t>0$. Put $L=\log x$. For each positive integer $m$ define
\begin{align}
 q_m&=4\pi m,& \alpha_m&=\frac1{8m},& \beta_m&=\frac12-\alpha_m,\label{eq:params}\\
 y_m(x)&=\frac{W(q_mx)}{q_m},&
 q_my_m+\log y_m&=L,\label{eq:saddle}\\
 \theta_m(x)&=\frac{q_my_m^2}{2}+y_m-\frac{\pi}{32m},&
 A_m(x)&=\frac{4x^{\alpha_m}y_m^{\beta_m}}{\sqrt{q_m}}.\label{eq:phase}
\end{align}
In a formula referring to one fixed mode we omit subscripts and write $q,\alpha,\beta,y,\theta,A$.

\begin{theorem}[All fixed orders of the discrepancy]\label{thm:main}
There are explicitly computable constants $d_{m,j}\in\mathbb Q(\ii)[\pi^{-1}]$ with $d_{m,0}=1$. For every fixed integer $K\ge1$,
\begin{equation}\label{eq:main}
 D(x)=-A_1(x)\Im\left\{e^{\ii\theta_1(x)}
          \sum_{j=0}^{K-1}\frac{d_{1,j}}{y_1(x)^j}\right\}
       +\OO_K\left(A_1(x)y_1(x)^{-K}\right).
\end{equation}
The same formula holds at positive integers with $D(n)$ replaced by $a(n)-2ne^n$. In particular,
\begin{align}
 d_{1,0}&=1,\notag\\
 d_{1,1}&=-\frac1{8\pi}+\frac{11\ii}{384},\label{eq:d123}\\
 d_{1,2}&=\frac3{128\pi^2}-\frac{2137}{294912}
                    +\frac{79\ii}{3072\pi}.\notag
\end{align}
The coefficient generator in Section~\ref{sec:coefficients} computes any prescribed fixed order by finite rational operations.
\end{theorem}

The leading formula reads
\begin{equation}\label{eq:leading}
 D(x)=-A_1(x)\sin\theta_1(x)+\OO(A_1(x)/y_1(x)).
\end{equation}
For example, the first corrected bracket is
\[
 D(x)=-A_1\left[\left(1-\frac1{8\pi y_1}\right)\sin\theta_1
                    +\frac{11}{384y_1}\cos\theta_1\right]
       +\OO(A_1y_1^{-2}).
\]
These are envelope estimates, including at every zero of the displayed sine. The error need not be small relative to the value of the oscillatory bracket.

\begin{corollary}[Sharp oscillation and the OEIS conjecture]\label{cor:sharp}
Let
\begin{equation}\label{eq:Cstar}
 C_*=4(4\pi)^{-7/8}.
\end{equation}
The set of subsequential limits of
\[
 \frac{a(n)-2ne^n}{n^{1/8}(\log n)^{3/8}}
\]
is exactly $[-C_*,C_*]$. In particular the limsup is $C_*$, the liminf is $-C_*$, and both positive and negative discrepancies are unbounded. Nevertheless $a(n)\sim2ne^n$. More strongly, for every fixed $r>0$,
\begin{equation}\label{eq:relative}
 \frac{a(n)}{2ne^n}=1+o(n^{-r}).
\end{equation}
Thus all ordinary inverse-power relative corrections vanish; the informative next scale is the oscillatory absolute error.
\end{corollary}

The fixed-order expansion is accompanied by a distinct exact-mode theorem. Define the absolutely convergent integral
\begin{equation}\label{eq:J}
 J_m(x)=\int_0^{\infty e^{\ii\pi/4}}
        \frac{2x^z}{\Gamma(z)}e^{2\pi\ii m z^2}\dd z.
\end{equation}
Since $1/\Gamma(z)$ is entire and vanishes at zero, this integral has no endpoint singularity. It equals the same integral over the positive real axis.

\begin{theorem}[Finite exact harmonic resolution]\label{thm:finite}
For every fixed integer $M\ge0$,
\begin{equation}\label{eq:finite}
 S(x)=M(x)-2N_1(x)+\frac12
           +2\Re\sum_{m=1}^M J_m(x)+R_{M+1}(x),
\end{equation}
where the sum is empty for $M=0$,
\begin{equation}\label{eq:N1}
 N_1(x)=-\int_0^\infty
       \frac{e^{-v}}{(\log v-\log x)^2+\pi^2}\dd v
       =\OO((\log x)^{-2}),
\end{equation}
and, for every fixed $\eps>0$,
\begin{equation}\label{eq:tail}
 R_{M+1}(x)=\OO_{M,\eps}
             \left(x^{1/(8(M+1))+\eps}\right).
\end{equation}
For every fixed $m\ge1$ and $K\ge1$ the mode itself has the expansion
\begin{align}
 2\Re J_m(x)
  ={}&-A_m(x)\Im\left\{e^{\ii\theta_m(x)}
                 \sum_{j=0}^{K-1}d_{m,j}y_m(x)^{-j}\right\}\notag\\
     &+\OO_{m,K}(A_m(x)y_m(x)^{-K})+\OO_m(L^{-2}).\label{eq:mode}
\end{align}
\end{theorem}

The parameters in \eqref{eq:tail} and \eqref{eq:mode} have different roles. The first controls the entire omitted tail after a fixed number of \emph{exact} modes. The second expands one fixed mode. Neither theorem asserts estimates uniform when $m$, $M$, or $K$ grows with $x$.

\section{Analytic preliminaries and the continuous contribution}\label{sec:prelim}
\subsection{Convergence and monotonicity}
For $x$ in any compact subset of the slit plane $\C\setminus(-\infty,0]$, use the principal logarithm in $x^u=e^{u\Log x}$. Real Stirling estimates give, for large $u$,
\[
 \frac{|x|^u}{\Gamma(1+u)}
 \le C\exp\{-u\log u+C_Ku\}.
\]
There are $\OO(u)$ indices $k$ with $u\le\sqrt k<u+1$. The series in \eqref{eq:model}, and the differentiated series on each compact set, therefore converge uniformly. Thus $S$ is holomorphic in that slit plane. On the positive axis it is strictly increasing, since each nonconstant summand increases. The same estimates, with $0<x\le1$, justify $S(x)\to1$ as $x\downarrow0$. Also $S(x)\ge1+x$, so $S(x)\to\infty$.

\subsection{The classical Volterra identity}
The Volterra function
\[
 \nu(x)=\int_0^\infty\frac{x^u}{\Gamma(1+u)}\dd u
\]
has the Ramanujan representation \cite[Theorem 3.1]{gm}
\begin{equation}\label{eq:ramanujan}
 \nu(x)=e^x-N_0(x),\qquad
 N_0(x)=\int_0^\infty
 \frac{e^{-xr}}{r((\log r)^2+\pi^2)}\dd r.
\end{equation}
The integral is convergent at zero because $1/(r(\log r)^2)$ is integrable there. For positive $x$, differentiation under the integral is justified locally by domination. Hence
\begin{align}
 I(x)=2x\nu'(x)&=2xe^x-2N_1(x),\label{eq:continuous}\\
 N_1(x):=xN_0'(x)&=-x\int_0^\infty
                 \frac{e^{-xr}}{(\log r)^2+\pi^2}\dd r.
\end{align}
The substitution $v=xr$ gives \eqref{eq:N1} and the elementary global bound $|N_1(x)|\le\pi^{-2}$.

For the sharper estimate, suppose $L=\log x>0$ and split the $v$ integral at $e^{L/2}$. On $0<v\le e^{L/2}$, $\log v-L\le-L/2$, and this part is at most $4L^{-2}$. The remaining part is at most $\pi^{-2}e^{-e^{L/2}}$. Thus
\begin{equation}\label{eq:contsmall}
 0<I(x)-2xe^x=-2N_1(x)=\OO(L^{-2}).
\end{equation}
No discrete oscillation is hidden in this continuous term. Its full inverse-logarithm asymptotics are part of the classical Volterra theory; see Garrappa--Mainardi and their discussion of Wyman--Wong \cite{gm}. The subsequent analysis addresses the difference between the integral and the square-root samples.

\section{Abel Plana at the square root endpoint}\label{sec:abel}
Set
\[
 f_x(z)=\frac{x^z}{\Gamma(1+z)},\qquad F_x(w)=f_x(\sqrt w),
\]
with the principal square root in the right half-plane. Although $F_x$ is not analytic at zero, its expansion there is
\begin{equation}\label{eq:endpoint}
 F_x(w)=1+(L+\gamma)\sqrt w+\OO_x(w),
\end{equation}
where $\gamma$ is Euler's constant.

\begin{lemma}[Endpoint Abel--Plana formula]\label{lem:abel}
For real $x>0$,
\begin{equation}\label{eq:abel}
 S(x)=I(x)+\frac12
 -2\int_0^\infty
       \frac{\Im F_x(\ii t)}{e^{2\pi t}-1}\dd t.
\end{equation}
The real integral on the right converges absolutely. In particular, it is the integral of the imaginary part, not an ordinary complex integral of $F_x(\ii t)/(e^{2\pi t}-1)$.
\end{lemma}

\begin{proof}
Apply the classical Abel--Plana formula \cite[Section 2.10]{dlmfap} to $w\mapsto F_x(w+\delta)$, $\delta>0$. This is analytic on a neighborhood of the closed right half-plane. Its positive real values decay faster than any power. On the right half-plane, complex Stirling gives the coarse bound
\[
 |F_x(w+\delta)|\le
 \exp\{C_x\sqrt{1+|w|}\log(2+|w|)\},
\]
for $\delta$ in a bounded interval. This is subexponential in $|w|$ and is dominated by the Abel--Plana kernel high on the vertical boundaries. To check the right edge also near real height zero, take rectangles with $\Re w=R+1/2$. When $|\Im w|\le R$, the square root has real part comparable with $\sqrt R$ and lies in a fixed narrow sector; Stirling gives $|F_x(w+\delta)|\le\exp(-c\sqrt R\log R)$ for large $R$, uniformly along that part. The summation kernel is bounded on the half-integer edge. When $|\Im w|\ge R$, its exponential decay dominates the displayed subexponential bound. Both right-edge integrals therefore vanish. Standard finite-contour truncation followed by an infinite limit is legitimate.

It remains to let $\delta\downarrow0$. Near zero, the difference of the two boundary values has the uniform estimate
\[
 |F_x(\delta+\ii t)-F_x(\delta-\ii t)|
       \le C_x\sqrt t\qquad(0<t\le1,\ 0<\delta\le1).
\]
To see this, $f_x'$ is bounded on a compact set containing the two square roots and the segment joining them; integrate the derivative along that segment and use
$|\sqrt{\delta+\ii t}-\sqrt{\delta-\ii t}|\le C\sqrt t$. Dividing by $e^{2\pi t}-1$ gives an integrable majorant $C_xt^{-1/2}$. The subexponential estimate supplies a common integrable majorant for $t\ge1$. Dominated convergence applies to the correction term. It also applies to the shifted real integral and sum by the gamma-decay estimate. Their endpoint value tends to $F_x(0)=1$.

For real $x$, the lower boundary value is the complex conjugate of the upper one. Thus the usual correction
\[
 \ii\int_0^\infty
 \frac{F_x(\ii t)-F_x(-\ii t)}{e^{2\pi t}-1}\dd t
\]
equals the last term in \eqref{eq:abel}. This completes the justification at the branch point.
\end{proof}

The distinction in the lemma is essential. Since $F_x(\ii t)=1+\OO_x(\sqrt t)$, the individual complex integral has a logarithmically divergent real part at zero. Removing that divergent part after an unjustified interchange can change an endpoint constant.

\section{Exact harmonics and the summed tail}\label{sec:tail}
\subsection{The exact kernel}
Write
\begin{equation}\label{eq:hK}
 h_x(z)=2z f_x(z)=\frac{2x^z}{\Gamma(z)},\qquad
 K_p(z)=\frac{e^{2\pi\ii pz^2}}{1-e^{2\pi\ii z^2}}\quad(p\ge1).
\end{equation}
On the ray $\arg z=\pi/4$, put $z=\sqrt{\ii t}$, so $\dd t=-2\ii z\dd z$. The correction in \eqref{eq:abel} becomes
\begin{equation}\label{eq:rayidentity}
 S(x)-I(x)-\frac12
 =2\lim_{\rho\downarrow0}\Re
    \int_{\rho e^{\ii\pi/4}}^{\infty e^{\ii\pi/4}}
           h_x(z)K_1(z)\dd z.
\end{equation}
The limit of the real part exists. Indeed, uniformly on a small angular sector,
\begin{equation}\label{eq:singular}
 h_x(z)=2z+\OO_x(z^2),\qquad
 h_x(z)K_p(z)=\frac{\ii}{\pi z}+\OO_x(1).
\end{equation}
A radial integral of the singular term is purely imaginary.

For every finite $M$ the elementary identity
\begin{equation}\label{eq:geometric}
 K_1(z)=\sum_{m=1}^M e^{2\pi\ii mz^2}+K_{M+1}(z)
\end{equation}
holds off the poles. Subtracting these finitely many absolutely convergent ray integrals proves \eqref{eq:finite}, with
\begin{equation}\label{eq:Rexact}
 R_{M+1}(x)=2\lim_{\rho\downarrow0}\Re
 \int_{\rho e^{\ii\pi/4}}^{\infty e^{\ii\pi/4}}
                 h_x(z)K_{M+1}(z)\dd z.
\end{equation}
No interchange of an infinite sum of singular endpoint integrals is involved.

\subsection{Deformation and the origin constant}
Fix $p=M+1$ and choose
\begin{equation}\label{eq:cchoice}
 \frac1{8p}<c<1.
\end{equation}
Deform the ray to the real segment $[0,c]$ followed by $c+\ii[0,\infty)$. The nonzero kernel poles satisfy $z^2\in\mathbb Z$ and lie on the real or imaginary axes. The real segment ends before the first positive real pole, and the vertical contour lies strictly to the right of the imaginary-axis poles. None is crossed.

The small counterclockwise arc $z=\rho e^{\ii\phi}$, $0\le\phi\le\pi/4$, has limiting integral
\begin{equation}\label{eq:arc}
 \int_0^{\pi/4}\frac{\ii}{\pi\rho e^{\ii\phi}}
                  \ii\rho e^{\ii\phi}\dd\phi=-\frac14.
\end{equation}
Moving this arc to the other side of the contour identity gives
\begin{align}
 \lim_{\rho\downarrow0}\Re\int_{\rho e^{\ii\pi/4}}^{\infty e^{\ii\pi/4}}
                 h_xK_p\dd z
 ={}&\frac14+\lim_{\rho\downarrow0}\Re\int_\rho^c h_x(t)K_p(t)\dd t\notag\\
    &+\Re\int_{c}^{c+\ii\infty}h_x(z)K_p(z)\dd z.\label{eq:deform}
\end{align}
The constant in $R_p$ is consequently $+1/2$, not $-1/2$ or zero.

For completeness the upper closing integrals vanish. Close horizontally from $c+\ii T$ to $T+\ii T$, where the latter point lies on the original ray. Uniform first-quadrant Stirling, for $c\le u\le T$, gives
\[
 |h_x(u+\ii T)|\le C_{c,x}T^{1/2}
       \exp\{\pi T/2+u(L+1-\log T)\}.
\]
Here one uses $T\le|u+\ii T|\le\sqrt2T$ and $\arg(u+\ii T)\le\pi/2$ in the logarithmic Stirling formula. The kernel is at most $C_ce^{-4\pi puT}$, since its denominator is uniformly separated from zero for large $T$. Once $\log T\ge L+1$, the entire closing integral is consequently bounded by
\[
 C_{c,x}T^{3/2}e^{-(4\pi pc-\pi/2)T}\longrightarrow0.
\]
This deformation concerns each fixed $x$; the uniform-in-$x$ estimate comes next from the resulting vertical contour.

On the real segment, \eqref{eq:geometric} gives the nonsingular real part
\begin{equation}\label{eq:realK}
 \Re K_p(t)=-\frac12-\sum_{m=1}^{p-1}\cos(2\pi mt^2),\qquad 0<t<c.
\end{equation}
Consequently \eqref{eq:deform} can be written entirely with ordinary convergent integrals:
\begin{align}
 R_p(x)={}&\frac12-\int_0^c h_x(t)\dd t
       -2\sum_{m=1}^{p-1}\int_0^c h_x(t)\cos(2\pi mt^2)\dd t\notag\\
       &+2\Re\int_0^\infty \ii h_x(c+\ii v)K_p(c+\ii v)\dd v.\label{eq:Rdeformed}
\end{align}
This also verifies the sign of the arc constant without leaving any divergent integral in the final formula.

\subsection{A bound for every omitted mode together}
For $x\ge1$, the horizontal integrals in \eqref{eq:Rdeformed} are $\OO_{c,p}(x^c)$. On the vertical line, complex Stirling gives
\begin{equation}\label{eq:hbound}
 |h_x(c+\ii v)|\le C_cx^c(1+v)^{1/2-c}e^{\pi v/2}.
\end{equation}
The denominator of $K_p$ is bounded away from zero on that line. At $v=0$, $c^2\notin\mathbb Z$; for $v>0$ the exponential in the denominator has modulus $e^{-4\pi cv}<1$, and it tends to zero. Compactness on a bounded initial interval and this last observation prove a uniform lower bound. Hence
\begin{equation}\label{eq:Kbound}
 |K_p(c+\ii v)|\le C_c e^{-4\pi pcv}.
\end{equation}
The decay margin $4\pi pc-\pi/2$ is positive by \eqref{eq:cchoice}. Absolute integration proves
\begin{equation}\label{eq:Rpower}
 R_p(x)=\OO_{c,p}(x^c).
\end{equation}
For a requested exponent $1/(8p)+\eps<1$, take that value as $c$; for a larger exponent choose a smaller admissible $c$. This proves \eqref{eq:tail}. Differentiating the deformed integrals any fixed number of times is also legitimate: the extra factors are polynomials in $c+\ii v$ times $x^{-r}$, still exponentially integrable. The constant term has derivative zero. Thus $R_p^{(r)}(x)=\OO_{c,p,r}(x^{c-r})$ for $r\ge1$.

\section{A full Gaussian line for each mode}\label{sec:gaussian}
\subsection{Rotation and completion}
Fix $m\ge1$, set $q=4\pi m$, and let $\omega=e^{\ii\pi/4}$. The positive-real representation
\[
 J_m(x)=\int_0^\infty\frac{2x^t}{\Gamma(t)}e^{\ii qt^2/2}\dd t
\]
is absolutely convergent by real gamma decay. Rotation to $\omega[0,\infty)$ is valid: away from the real axis the quadratic exponential decays, while near that axis the reciprocal-gamma factor supplies the needed decay on the large arc. At zero the integrand is $\OO(z)$, so there is no arc contribution.

Complete the ray to the whole oriented line by defining
\begin{align}
 F_m(x)&=\int_{\omega\R}\frac{2e^{Lz}}{\Gamma(z)}e^{\ii qz^2/2}\dd z,\label{eq:F}\\
 K_m^-(x)&=\int_{\omega(-\infty,0]}\frac{2e^{Lz}}{\Gamma(z)}e^{\ii qz^2/2}\dd z.
\end{align}
Then $J_m=F_m-K_m^-$. The notation $K_m^-$ is unrelated to the tail kernel $K_p(z)$.

\begin{lemma}[The added negative ray]\label{lem:negative}
As $L\to+\infty$, $K_m^-(x)=\OO_m(L^{-2})$.
\end{lemma}
\begin{proof}
The reciprocal gamma function is entire of order one and satisfies a global bound
\begin{equation}\label{eq:globalgamma}
 |1/\Gamma(z)|\le\exp\{C(1+|z|)\log(2+|z|)\}.
\end{equation}
This follows, for example, from sectorial Stirling on the right and upper/lower sectors together with the reflection identity on the remaining sector. Its simple zero at zero and the domination of $s\log(2+s)$ by every positive multiple of $s^2$ imply, for fixed $q$, a finite constant $B_q$ such that
\[
 |1/\Gamma(-\omega s)|\le B_qs e^{qs^2/4}\qquad(s\ge0).
\]
Along that ray $|e^{Lz}|=e^{-Ls/\sqrt2}$ and
$|e^{\ii qz^2/2}|=e^{-qs^2/2}$. Therefore
\[
 |K_m^-|\le2B_q\int_0^\infty
       s e^{-Ls/\sqrt2-qs^2/4}\dd s
       \le\frac{4B_q}{L^2}.
\]
\end{proof}
The added ray is negligible relative to $x^{\alpha_m}y_m^{\beta_m-K}$ at every fixed order. Completion avoids artificial endpoint errors that can arise when a short real segment and an unbounded vertical half-line are expanded separately.

\subsection{Shifting to the saddle}
Let $z_0=\alpha+\ii y$, with the parameters in \eqref{eq:params}--\eqref{eq:saddle}. The integrand of $F_m$ is entire, so the line may be shifted from $\omega\R$ to $z_0+\omega\R$. The quadratic Gaussian dominates \eqref{eq:globalgamma} on the connecting ends, proving that these vanish.

Writing $z=z_0+\omega r$, $r\in\R$, gives the exact identity
\begin{equation}\label{eq:cancellation}
 Lz+\frac{\ii qz^2}{2}
 =Lz_0+\frac{\ii qz_0^2}{2}
        +\omega\log(\ii y)r-\frac{qr^2}{2}.
\end{equation}
Indeed, $L-qy=\log y$ and $q\alpha=\pi/2$. The linear term is precisely the leading linear term of $\log\Gamma(z)$ near $z_0$.

\begin{lemma}[Global localization]\label{lem:localization}
After division by its value at $z_0$, the modulus of the shifted integrand is bounded by a fixed integrable Gaussian for $|r|\le\eta y$, where $\eta>0$ is sufficiently small and fixed. The integral over $|r|>\eta y$ is exponentially small in $y^2$ relative to $x^\alpha y^\beta$. Consequently, for any fixed $0<\delta<1/3$, restricting the integral to $|r|\le y^\delta$ loses an error smaller than $x^\alpha y^{\beta-K}$ for every fixed $K$.
\end{lemma}
\begin{proof}
Choose $\eta$ so that $z_0+\omega r$ stays in a fixed closed sector away from the negative real axis and has modulus comparable with $y$ when $|r|\le\eta y$. Uniform Stirling and its differentiated form give
\begin{equation}\label{eq:loggammadiff}
 \log\Gamma(z_0+\omega r)-\log\Gamma(z_0)
   =\omega r\log(\ii y)+\OO_m((|r|+r^2)/y).
\end{equation}
For clarity, the derivative satisfies
$\psi(z_0)-\log(\ii y)=\OO_m(y^{-1})$, and $\psi'(z)=\OO_m(y^{-1})$ throughout this segment; integrating these bounds proves \eqref{eq:loggammadiff}. Combining it with \eqref{eq:cancellation} bounds the normalized modulus by
\[
 C_m\exp\left\{-\frac{qr^2}{2}+C_m\frac{|r|+r^2}{y}\right\}
 \le C'_m e^{-qr^2/4}
\]
for all sufficiently large $y$.

For $|r|\ge\eta y$ we do not use sectorial Stirling on a contour that may cross the negative real axis. Use \eqref{eq:globalgamma} instead. The real part of the linear term in \eqref{eq:cancellation} is $\OO_m(|r|\log y)$. The value at $z_0$ has size comparable with $x^\alpha y^\beta$, by Stirling at $z_0$. Up to this scale the logarithm of a global majorant is
\begin{equation}\label{eq:farglobal}
 -\frac{qr^2}{2}
 +\OO_m\bigl(|r|\log y+(y+|r|)\log(2+y+|r|)\bigr).
\end{equation}
Uniformly for $|r|\ge\eta y$, the error divided by $r^2$ tends to zero as $y\to\infty$. Thus \eqref{eq:farglobal} is at most $-qr^2/4$ there for large $y$. Its integral is $\OO_m(e^{-c_my^2})$. On $y^\delta<|r|\le\eta y$ the previous local Gaussian gives $\OO_m(e^{-c_my^{2\delta}})$. Both estimates dominate any prescribed inverse power of $y$.
\end{proof}

This global step is what permits an all-orders conclusion. A formal stationary point calculation alone would not exclude other contour portions of comparable size.

\section{Every fixed order and a finite coefficient generator}\label{sec:coefficients}
\subsection{Uniform Stirling expansion near the saddle}
For fixed real $\alpha$ and every fixed integer $N\ge1$, the shifted Stirling expansion is
\begin{align}
 \log\Gamma(\alpha+\ii t)
 ={}&(\alpha+\ii t-\tfrac12)\log(\ii t)-\ii t
       +\tfrac12\log(2\pi)\notag\\
 &+\sum_{r=1}^{N-1}
 \frac{(-1)^{r+1}B_{r+1}(\alpha)}{r(r+1)(\ii t)^r}
       +\OO_{\alpha,N}(|t|^{-N}),\label{eq:shiftedstirling}
\end{align}
uniformly in the fixed closed sectors used locally above \cite[Section 5.11]{dlmfgamma}. Here $B_j(\alpha)$ is the Bernoulli polynomial. The logarithm is the analytic branch inherited from that local sector. The same uniform control applies on a slightly larger sector, so derivatives needed for Taylor estimates follow by Cauchy's estimate.

In the vertical coordinate $z=\alpha+\ii t$, the cancellation $q\alpha=\pi/2$ yields the local asymptotic integrand
\begin{align}
 \ii\frac{2x^{\alpha+\ii t}}{\Gamma(\alpha+\ii t)}
                  e^{\ii q(\alpha+\ii t)^2/2}
 \sim{}&\ii\sqrt{\frac2\pi}\,x^\alpha t^\beta
     e^{\ii(\Phi(t)+\pi/4-\pi/(32m))}\notag\\
 &\times\exp\left\{-\sum_{r\ge1}
  \frac{(-1)^{r+1}B_{r+1}(\alpha)}{r(r+1)(\ii t)^r}\right\},\label{eq:verticalformal}\\
 \Phi(t)&=t(L-\log t+1)-\frac{qt^2}{2}.\label{eq:Phi}
\end{align}
The infinite expression in \eqref{eq:verticalformal} is a \emph{formal asymptotic series}. At any use it is truncated at a fixed order and accompanied by the remainder in \eqref{eq:shiftedstirling}. No convergence of that infinite Stirling series is asserted.

The saddle equation gives $\Phi'(y)=0$ and
\begin{equation}\label{eq:phiexpand}
 \Phi(y)=\frac{qy^2}{2}+y,\qquad
 \Phi(y+v)-\Phi(y)
 =-\frac{qv^2}{2}
 +\sum_{r\ge2}\frac{(-1)^{r-1}v^r}{r(r-1)y^{r-1}}.
\end{equation}
The latter Taylor expansion converges locally for $|v|<y$, though only finitely many terms are required. On the shifted line, $v=e^{-\ii\pi/4}r$; the leading quadratic factor is a decaying Gaussian.

\subsection{The Bernoulli and Gaussian algebra}
Define a linear functional on polynomials by
\begin{equation}\label{eq:G}
 \G_q(v^{2r})=(-\ii/q)^r(2r-1)!!,\qquad
 \G_q(v^{2r+1})=0,\qquad \G_q(1)=1.
\end{equation}
These are the normalized moments of the Gaussian along $e^{-\ii\pi/4}\R$. Let $\eps$ be a formal variable and define
\begin{align}
 H(\eps,v)={}&(1+\eps v)^\beta\notag\\[-4pt]
 &\times\exp\left\{-\sum_{r\ge1}
 \frac{(-1)^{r+1}B_{r+1}(\alpha)}{r(r+1)}
       \left(\frac{\eps}{\ii(1+\eps v)}\right)^r
 +\ii\sum_{r\ge2}\frac{(-1)^{r-1}\eps^{r-1}v^r}{r(r-1)}\right\}.\label{eq:H}
\end{align}
Then the desired constants are
\begin{equation}\label{eq:dgen}
 d_{m,j}=[\eps^j]\G_q(H(\eps,v)).
\end{equation}
This is well-defined formal algebra: for a fixed coefficient $\eps^j$, only finitely many summands and polynomial degrees are involved.

For an explicit finite recurrence, write
\[
 C_r=\frac{(-1)^{r+1}B_{r+1}(\alpha)}{r(r+1)},\qquad
 \log H=\sum_{j\ge1}\ell_j(v)\eps^j.
\]
A direct expansion gives
\begin{align}
 \ell_j(v)={}&\frac{\beta(-1)^{j+1}}j v^j
 -\sum_{r=1}^j C_r\ii^{-r}(-1)^{j-r}
                 \binom{j-1}{j-r}v^{j-r}
 +\frac{\ii(-1)^j}{j(j+1)}v^{j+1}.\label{eq:ell}
\end{align}
If $H=\sum_{j\ge0}H_j(v)\eps^j$, the exponential identity implies
\begin{equation}\label{eq:Hrec}
 H_0=1,\qquad H_j=\frac1j\sum_{r=1}^j r\ell_rH_{j-r},\qquad
 d_{m,j}=\G_q(H_j).
\end{equation}
Bernoulli numbers may be generated from $B_0=1$ and
\[
 B_n=-\frac1{n+1}\sum_{k=0}^{n-1}\binom{n+1}{k}B_k,
 \qquad B_n(\alpha)=\sum_{k=0}^n\binom nkB_k\alpha^{n-k}.
\]
Together these formulas use only rational complex arithmetic, with each coefficient stored as a finite polynomial in $\pi^{-1}$. In particular there is no need to evaluate $\pi$ to compute a coefficient exactly.

\subsection{Why the formal calculation is an asymptotic theorem}
The base Gaussian integral is
\[
 \int_{e^{-\ii\pi/4}\R}e^{-\ii qv^2/2}\dd v
       =e^{-\ii\pi/4}\sqrt{2\pi/q}.
\]
Combining it with the constant factors in \eqref{eq:verticalformal} gives
\begin{equation}\label{eq:Fexpansion}
 F_m(x)=\frac{2x^\alpha y^\beta}{\sqrt q}
          e^{\ii(\theta+\pi/2)}
       \left(\sum_{j=0}^{K-1}d_{m,j}y^{-j}
                        +\OO_{m,K}(y^{-K})\right).
\end{equation}
We give the remainder justification rather than treating Gaussian moment manipulation as a proof.

By Lemma~\ref{lem:localization}, restrict first to $|r|\le y^\delta$, for a fixed $0<\delta<1/3$. On this portion $|v|/y\to0$, and the finite Stirling remainder in \eqref{eq:shiftedstirling} is $\OO_{m,N}(y^{-N})$, uniformly. Expand the algebraic amplitude, the finite inverse-gamma correction, and the phase correction in powers of $1/y$ to order $K-1$, using a Stirling truncation with $N\ge K$ and, if necessary, a larger fixed $N$ for intermediate Taylor remainders. The logarithm of the non-Gaussian multiplier is $\OO_m((1+|r|+r^2)/y)$ on this cutoff. In particular it is uniformly bounded and tends to zero there.

Taylor's theorem for the exponential and the rational/analytic factors now bounds the difference between the exact multiplier and its first $K$ formal coefficients by
\[
 C_{m,K}y^{-K}P_{m,K}(|r|)
       \exp\{C_m(1+|r|+r^2)/y\},
\]
where $P_{m,K}$ is a fixed polynomial with nonnegative coefficients. This follows equally from the integral remainder of each finite Taylor expansion; derivatives of the analytic factors on $|v|/y\le1/2$ are polynomially bounded in $|r|$. The dominating Gaussian from Lemma~\ref{lem:localization} absorbs the last exponential for large $y$. Its polynomial moments are finite, so the integrated remainder is $\OO_{m,K}(y^{-K})$ after normalization. Extending each finitely many polynomial--Gaussian integral from the cutoff to the full line changes it by an exponentially small error. Their values are exactly \eqref{eq:G}. This proves \eqref{eq:Fexpansion} with a fixed-order remainder.

Subtracting the negative ray from Lemma~\ref{lem:negative} and taking twice the real part gives \eqref{eq:mode}. The two uses of a contour are distinct: the vertical expression is a convenient coefficient formula; the actual estimates take place on a globally localized, decaying Gaussian line.

\subsection{Independent low order expansion}
For $m=1$, $\alpha=1/8$ and $\beta=3/8$. The first coefficient follows immediately from
\[
 H_1(v)=\beta v+\frac{\ii B_2(\alpha)}2-\frac{\ii v^2}2,
 \qquad d_{1,1}=-\frac1{2q}+\frac{\ii B_2(1/8)}2.
\]
Since $B_2(1/8)=11/192$ and $q=4\pi$, this is \eqref{eq:d123}.

As an independent check of the next order, retain the exact local quadratic curvature $Q=q+1/y$. At the saddle the reciprocal-gamma correction is
\[
 1+\frac{11\ii}{384y}-\frac{2137}{294912y^2}+\OO(y^{-3}).
\]
The second-derivative amplitude correction to stationary phase contributes
\[
 -\frac{\ii\beta(\beta-1)}{2qy^2}
       =\frac{15\ii}{512\pi y^2}.
\]
Terms from the cubic phase begin at higher order here: its coefficient is $\OO(y^{-2})$, its odd base Gaussian moment vanishes, and a nonzero contribution requires an additional amplitude derivative or phase term. Thus an equivalent three-term form is
\begin{align}
 2\Re J_1(x)=-\frac{4x^{1/8}y^{3/8}}{\sqrt Q}
 \left[\sin\theta+\frac{11}{384y}\cos\theta\right.\qquad\notag\\[-3pt]
 \left.\hspace{26mm}+\frac1{y^2}
 \left(-\frac{2137}{294912}\sin\theta
             +\frac{15}{512\pi}\cos\theta\right)
 +\OO(y^{-3})\right]+\OO(L^{-2}).\label{eq:curvatureform}
\end{align}
Expanding
\[
 \sqrt{q/Q}=1-\frac1{2qy}+\frac3{8q^2y^2}+\OO(y^{-3})
\]
recovers $d_{1,2}$ in \eqref{eq:d123}. This route checks the rational, real-curvature and imaginary-phase pieces separately.

\section{From one mode to sharp integer oscillation}\label{sec:oscillation}
\subsection{Proof of the main expansion}
Take $M=1$ in Theorem~\ref{thm:finite}, and choose a fixed $c$ with $1/16<c<1/8$. The exact remainder is $\OO(x^c)$. For every fixed $K$,
\[
 x^c+1+L^{-2}=o\bigl(x^{1/8}y_1^{3/8-K}\bigr),
\]
because $y_1\sim L/(4\pi)$. Combining \eqref{eq:continuous}, \eqref{eq:finite} and \eqref{eq:mode} proves \eqref{eq:main}. The floor changes the result by a quantity in $(-1,0]$, also absorbed by that error. This proves Theorem~\ref{thm:main}.

The equation $qy+\log y=L$ shows $y\sim L/q$ and gives, if desired, the familiar Lambert expansion
\[
 y=\frac1q\left[L-\log L+\log q
                +\OO\!\left(\frac{\log L}{L}\right)\right].
\]
Keeping $W(qx)$ exact is cleaner: a small phase error can matter more than a relative amplitude error. In particular,
\begin{equation}\label{eq:envelope}
 A_1(x)\sim C_*x^{1/8}(\log x)^{3/8}.
\end{equation}
The claimed leading equivalent and \eqref{eq:relative} follow since this envelope is exponentially smaller than $M(x)$.

\subsection{All subsequential limits}
Implicit differentiation of \eqref{eq:saddle} gives
\[
 \frac{\dd y_m}{\dd L}=\frac1{q_m+1/y_m},\qquad
 \frac{\dd\theta_m}{\dd L}=y_m,\qquad
 \theta_m'(x)=\frac{y_m(x)}x.
\]
Thus $\theta_1$ is increasing and unbounded, while $\theta_1'(x)\to0$. Fix any $s\in[-1,1]$ and choose $\phi$ with $-\sin\phi=s$. For every sufficiently large integer $j$, let $x_j$ be the unique real solution of
\[
 \theta_1(x_j)=\phi+2\pi j,
\]
and choose an integer $n_j$ nearest to $x_j$. By the mean value theorem,
\[
 |\theta_1(n_j)-\theta_1(x_j)|
 \le\tfrac12\sup_{|t-x_j|\le1/2}\theta_1'(t)\longrightarrow0.
\]
The sequence $n_j$ tends to infinity and has a strictly increasing subsequence. Dividing \eqref{eq:leading} by $A_1(n_j)$ shows that the normalized discrepancy tends to $s$. Conversely, the same formula confines all subsequential limits to $[-1,1]$. Rescaling by \eqref{eq:envelope} proves Corollary~\ref{cor:sharp}.

In particular, the discrepancy is not $\OO(n^{1/8-\delta})$ for any $\delta>0$, and $a(n)$ is not eventually $\lfloor2ne^n\rfloor$ or any uniformly bounded rounding modification of $2ne^n$. No equidistribution assertion is required for these conclusions.

\section{Complex sector control}\label{sec:sector}
A real-axis asymptotic estimate alone does not justify a convergent analytic inverse. We therefore prove the relevant complex bound for the exact discrepancy separately.

\begin{theorem}[Analytic sector bound]\label{thm:sector}
For every $a>1/8$, there are $\delta>0$, $C_a>0$ and $R_a>0$ such that $D$ is holomorphic and
\begin{equation}\label{eq:sectorbound}
 |D(z)|\le C_a|z|^a
 \qquad(|z|\ge R_a,\ |\arg z|<\delta).
\end{equation}
For every fixed integer $r\ge0$ this implies
\begin{equation}\label{eq:derivativebound}
 D^{(r)}(x)=\OO_{a,r}(x^{a-r})\qquad(x\to+\infty).
\end{equation}
\end{theorem}

\begin{proof}
Choose $c$ with $1/8<c<\min(a,1)$ and put $\mu=4\pi c-\pi/2>0$. Fix $0<\delta<\min(\mu/2,\pi/3)$. The defining series is holomorphic on the slit plane by Section~\ref{sec:prelim}. We construct a uniformly bounded continuation of the $M=0$ exact contour formula.

For a complex parameter $z$ write $h_z(\zeta)=2z^\zeta/\Gamma(\zeta)$. Pair the upper contour in \eqref{eq:deform} with its reflected lower contour. On the real horizontal segment their kernels sum to
\[
 K_1(t)+\overline{K_1(t)}=-1,
\]
so the divergent imaginary endpoint parts cancel exactly. The paired horizontal integral is therefore $-\int_0^c h_z(t)\dd t$, an analytic function bounded by $C_c|z|^c$ for $|z|\ge1$. The two small-arc constants add to $1/2$ in $R_1$.

More explicitly, if $U(z)$ denotes the upper vertical integral, its partner is $U^*(z)=\overline{U(\bar z)}$. This is holomorphic in $z$, not an operation of complex conjugation left inside the unknown function. It is equivalently the reflected, oppositely oriented lower contour integral. Their absolute estimates use
\[
 |z^{c+\ii v}|=|z|^ce^{-v\arg z},\qquad
 |z^{c-\ii v}|=|z|^ce^{v\arg z}.
\]
The reciprocal gamma and kernel estimates \eqref{eq:hbound}--\eqref{eq:Kbound} leave $e^{-\mu v}$. Both integrals are bounded uniformly in the sector by
\[
 C_c|z|^c\int_0^\infty
         (1+v)^{1/2-c}e^{-(\mu-\delta)v}\dd v
       =\OO_c(|z|^c).
\]
The same integrable majorant on compact subsets proves analyticity of these contour integrals.

The continuous correction has the analytic representation
\[
 I(z)-2ze^z=2z\int_0^\infty
            \frac{e^{-zr}}{(\log r)^2+\pi^2}\dd r,
\]
valid initially for positive real $z$ and analytically continued to $\Re z>0$. Its modulus is at most $2|z|/(\pi^2\Re z)$, which is uniformly bounded in our sector. Adding this term, the original Abel--Plana endpoint $1/2$, and the paired contour formula yields a holomorphic function equal to $D$ on the positive axis. The identity theorem identifies it with $D$ throughout the sector. The preceding estimates give $\OO(|z|^c)$, hence \eqref{eq:sectorbound}.

For large real $x$, a disk of radius $\kappa x$ about $x$ lies in a slightly smaller sector, with a fixed $\kappa>0$. Cauchy's derivative estimate on that disk proves \eqref{eq:derivativebound}.
\end{proof}

The same pairing could be used for higher exact-mode remainders. Only the $M=0$ case is needed below. In particular \eqref{eq:sectorbound} gives a uniform bound $|D(w+t)|\le C_aw^a$ on each sufficiently small fixed complex disk $|t|\le\rho$ about every large positive $w$.

\section{The continuous inverse and its convergent expansion}\label{sec:inverse}
Since $S$ is continuous and strictly increasing from $1$ to infinity, its real inverse
\[
 X(Y)=S^{-1}(Y)\qquad(Y>1)
\]
is well-defined. Put
\begin{equation}\label{eq:w}
 w=W(Y/2),\qquad M(w)=Y,\qquad M'(w)=Y(1+1/w).
\end{equation}
We first state an exact-discrepancy inverse expansion. Its convergence concerns powers of the exact analytic function $D$, not the formal Stirling series or an infinite harmonic asymptotic sum.

\begin{theorem}[Convergent Lagrange inverse]\label{thm:inverse}
Let
\begin{equation}\label{eq:EQ}
 E_w(t)=e^t(1+t/w)-1,\qquad Q_w(t)=\frac{t}{E_w(t)},
 \qquad Q_w(0)=\frac1{1+1/w}.
\end{equation}
For all sufficiently large $Y$, the following series converges absolutely to the real inverse displacement:
\begin{equation}\label{eq:lagrange}
 X(Y)-w=\sum_{r=1}^\infty\frac{(-1)^r}{r!Y^r}
  \left.\frac{\dd^{r-1}}{\dd t^{r-1}}
      \bigl[D(w+t)^rQ_w(t)^r\bigr]\right|_{t=0}.
\end{equation}
For each fixed $a>1/8$ and integer $R\ge1$, the remainder after $R$ terms is
\begin{equation}\label{eq:invrem}
 \OO_{a,R}\left((w^a/Y)^{R+1}\right).
\end{equation}
\end{theorem}

\begin{proof}
Take a fixed radius $\rho=1/4$. On $|t|\le\rho$,
\[
 E_w(t)/t\longrightarrow(e^t-1)/t
\]
uniformly as $w\to\infty$, with the removable values understood. The limiting function has no zero on this disk, because the nonzero zeros of $e^t-1$ are $2\pi\ii k$. Therefore $Q_w$ is analytic and uniformly bounded on a neighborhood of this closed disk for large $w$.

The exact equation $S(w+t)=Y$ is
\[
 YE_w(t)+D(w+t)=0,
 \qquad t=-Y^{-1}D(w+t)Q_w(t).
\]
By Theorem~\ref{thm:sector}, $\sup_{|t|\le\rho}|D(w+t)Q_w(t)|\le C_aw^a$. Introduce a parameter $\lambda$ and apply Rouch\'e's theorem to
$t+\lambda D(w+t)Q_w(t)$ on $|t|=\rho$. When $|\lambda|C_aw^a<\rho$, it has exactly one root, counted with multiplicity, inside the disk. The root is simple, is analytic in $\lambda$ near zero, and its local power series is the classical Lagrange--B\"urmann formula, namely \eqref{eq:lagrange} with $Y^{-1}$ replaced by $\lambda$. This is an application of standard analytic inversion, not a new inversion principle.

Cauchy's estimate bounds its $r$th term at $\lambda=1/Y$ by
\begin{align*}
 \frac1{r!Y^r}
 \left|\left.\frac{\dd^{r-1}}{\dd t^{r-1}}
           [D(w+t)^rQ_w(t)^r]\right|_{t=0}\right|
 \le \frac{\rho}{r}
           \left(\frac{C_aw^a}{\rho Y}\right)^r.
\end{align*}
The right side is summable for large $w$, proving absolute convergence and the fixed-$R$ geometric tail estimate \eqref{eq:invrem}. At real $\lambda=1/Y$, the unique root is real by conjugation symmetry. Its equation is exactly $S(w+t)=Y$, so strict monotonicity identifies it with $X(Y)-w$.
\end{proof}

For transparency the first two exact-discrepancy terms are
\begin{align}
 X(Y)-w={}&-\frac{D(w)}{Yb_w}\notag\\
 &+\frac1{Y^2}\left(
       \frac{D(w)D'(w)}{b_w^2}
       -\frac{D(w)^2c_w}{b_w^3}\right)
       +\OO_a((w^a/Y)^3),\label{eq:inv2}\\
 b_w&=1+1/w,\qquad c_w=1/2+1/w.\notag
\end{align}
Indeed $Q_w(0)=b_w^{-1}$ and $Q_w'(0)=-c_wb_w^{-2}$. Higher coefficients follow from the same finite differentiation rule once the exact analytic $D$ is supplied.

\begin{corollary}[Oscillatory inverse displacement]\label{cor:invosc}
For every fixed $a>1/8$,
\begin{equation}\label{eq:invfirst}
 X(Y)=w-\frac{D(w)}{Y(1+1/w)}
                  +\OO_a(w^{2a}/Y^2).
\end{equation}
For every fixed integer $K\ge1$,
\begin{align}
 X(Y)={}&w+\frac{A_1(w)}{Y(1+1/w)}
       \Im\left\{e^{\ii\theta_1(w)}
                   \sum_{j=0}^{K-1}d_{1,j}y_1(w)^{-j}\right\}\notag\\
 &+\OO_K\left(\frac{A_1(w)y_1(w)^{-K}}Y\right).\label{eq:invosc}
\end{align}
\end{corollary}
\begin{proof}
The $r=1$ term and the remainder of Theorem~\ref{thm:inverse} give \eqref{eq:invfirst}. Substitute Theorem~\ref{thm:main}. For a fixed choice of $a>1/8$, the quadratic inverse error $w^{2a}/Y^2$ is smaller than $A_1(w)y_1(w)^{-K}/Y$ for every fixed $K$, since $Y=2we^w$. This proves \eqref{eq:invosc}.
\end{proof}

Thus the first inverse-displacement scale is exponentially small in $w$, but has all fixed logarithmic refinements. The convergent series \eqref{eq:lagrange} also resolves further exact-discrepancy powers. Substituting an asymptotic approximation for $D$ into it requires carrying its own error; it does not turn the fixed-order Stirling expansion into a convergent formula.

\section{Exact integer thresholds and failure of a naive ceiling}\label{sec:thresholds}
For an integer $Y\ge2$ define
\begin{equation}\label{eq:N}
 N(Y)=\min\{n\ge0:a(n)\ge Y\}.
\end{equation}
The equivalence $\lfloor S(n)\rfloor\ge Y\iff S(n)\ge Y$ gives the exact identity
\begin{equation}\label{eq:ceilexact}
 N(Y)=\lceil X(Y)\rceil.
\end{equation}
For a real threshold $Y>1$, first replace $Y$ by $\lceil Y\rceil$. If $Y\le1$, the threshold index is zero.

A proved pointwise error $|\widehat X(Y)-X(Y)|\le\eta(Y)$ yields the reliable enclosure
\begin{equation}\label{eq:ceilenclosure}
 \lceil\widehat X(Y)-\eta(Y)\rceil
       \le N(Y)\le
 \lceil\widehat X(Y)+\eta(Y)\rceil.
\end{equation}
When both ceilings agree, they determine $N(Y)$ exactly. Our asymptotic big-$O$ results do not supply effective numerical constants or an artificial uniform gap from integers; they do not by themselves provide certified finite-$Y$ values in this enclosure.

\begin{theorem}[Both directions of ceiling failure]\label{thm:failure}
For large integer $Y$, $N(Y)$ and $\lceil W(Y/2)\rceil$ differ by at most one. Equality nevertheless fails infinitely often in each direction. More precisely, for infinitely many integers $n$,
\begin{align}
 Y_n^+&=\lfloor M(n)\rfloor+1,&
 N(Y_n^+)&=n,& \lceil W(Y_n^+/2)\rceil&=n+1,\label{eq:failplus}\\
 Y_n^-&=\lfloor M(n)\rfloor-1,&
 N(Y_n^-)&=n+1,& \lceil W(Y_n^-/2)\rceil&=n.\label{eq:failminus}
\end{align}
\end{theorem}
\begin{proof}
Equation~\eqref{eq:invfirst} implies $X(Y)-W(Y/2)\to0$. For large $Y$ the difference has modulus less than one, and the ceilings of two real numbers at distance less than one differ by at most one.

Corollary~\ref{cor:sharp} supplies infinitely many $n$ with $D(n)>2$, and infinitely many with $D(n)<-2$. For the former set, define $Y_n^+$ as in \eqref{eq:failplus}. Then
\[
 M(n)<Y_n^+\le M(n)+1<S(n).
\]
The neighboring main-term gaps satisfy
\[
 M(n)-M(n-1)\asymp ne^n,\qquad M(n+1)-M(n)\asymp ne^n,
\]
whereas $D(n)$ is polynomially bounded. Thus, for large such $n$, $S(n-1)<Y_n^+\le S(n)$ and $M(n)<Y_n^+<M(n+1)$. Strict monotonicity gives the two equalities in \eqref{eq:failplus}.

For $D(n)<-2$, the number $Y_n^-\in(M(n)-2,M(n)-1]$ satisfies
\[
 S(n)<Y_n^-<M(n).
\]
The same gap estimate gives $S(n)<Y_n^-\le S(n+1)$ and $M(n-1)<Y_n^-<M(n)$. This proves \eqref{eq:failminus}. The thresholds are positive and tend to infinity along either subsequence.
\end{proof}

A small real inverse error can therefore change a carefully placed integer threshold. The sign-changing, unbounded direct discrepancy is precisely what produces both families of counterexamples.

\section{Further structural consequences}\label{sec:consequences}
\begin{proposition}[Increase and eventual strict log concavity]\label{prop:logconcavity}
The integer sequence $a(n)$ is strictly increasing for every $n\ge0$ and eventually satisfies
\[
 a(n)^2>a(n-1)a(n+1).
\]
\end{proposition}
\begin{proof}
The $k=1$ summand gives $S(n+1)-S(n)\ge1$, and the other nonconstant positive summands make the inequality strict. Hence the floors increase by at least one. The same statement holds from $n=0$ with $S(0)=1$.

Fix $a>1/8$; the discrepancy theorem gives $a(n)=M(n)+\OO(n^a)$. A direct calculation yields
\[
 M(n)^2-M(n-1)M(n+1)
 =4e^{2n}\bigl(n^2-(n-1)(n+1)\bigr)=4e^{2n}.
\]
Replacing each $M$ by $a$ changes this expression by
$\OO(e^nn^{a+1}+n^{2a})$, which is smaller than $4e^{2n}$ for all sufficiently large $n$. This proves strict log concavity.
\end{proof}

There is also a transparent hierarchy of scales. The $m$th mode has envelope
\[
 A_m(x)\asymp_m x^{1/(8m)}(\log x)^{1/2-1/(8m)}.
\]
Every fixed lower mode eventually dominates every fixed logarithmic correction to a higher-index mode in absolute envelope, whereas every fixed logarithmic correction of the first mode dominates the combined exact remainder after that mode, by the fixed-cutoff tail theorem. This is an envelope comparison, not a pointwise ordering of oscillatory values near their zeros.

\section{Verification and reproducibility}\label{sec:verification}
\subsection{What the exact checks establish}
The mandatory verification uses Python's standard library, integer arithmetic and rational complex pairs. It implements \eqref{eq:ell}--\eqref{eq:Hrec} and the Gaussian moments \eqref{eq:G}, generates $d_{m,0},\ldots,d_{m,6}$ for $m=1,2$, and checks the displayed coefficients exactly. Input guards reject unsupported or malformed parameters instead of silently computing a different model. Validation uses explicit exceptions, not assertions that disappear under optimized Python.

The accompanying corruption tests deliberately alter coefficient data, tail records and integrity metadata and require rejection. Running under both ordinary Python and \texttt{python -O} verifies that optimization does not remove these checks. These tests validate finite algebra and package behavior; the analytic estimates in this report remain mathematical proofs, not consequences of a finite numerical test.

\subsection{An exact bound for a positive omitted tail}
For a positive integer $n$ and an integer $T>n$, let
\[
 g_n(u)=\frac{n^u}{\Gamma(1+u)},\qquad c=\frac{T-n}{T}>0.
\]
Log convexity of $\Gamma$ gives, for $u\ge T$,
\[
 \log\Gamma(1+u)-\log\Gamma(1+T)
    \ge(u-T)\bigl(\log\Gamma(1+T)-\log\Gamma(T)\bigr)
    =(u-T)\log T.
\]
The same convexity shows $g_n$ is decreasing on $[T,\infty)$, since its logarithmic derivative is at most $\log(n/T)<0$. Also $\log(n/T)=\log(1-c)\le-c$, so
\[
 g_n(u)\le\frac{n^T}{T!}e^{-c(u-T)}.
\]
The integral test for the decreasing function $t\mapsto g_n(\sqrt t)$ therefore yields the exact bound
\begin{equation}\label{eq:numerictail}
 0<\sum_{k>T^2}\frac{n^{\sqrt k}}{\Gamma(1+\sqrt k)}
 \le2\frac{n^T}{T!}\left(\frac{T}{c}+\frac1{c^2}\right).
\end{equation}
Every quantity on the right is rational. For $n=1,\ldots,34$ and $T=4n+60$, exact comparison shows this bound is below $10^{-35}$. At $n=0$ the series is exactly one. The mandatory code checks these rational inequalities; it does not need an implementation of gamma, logarithms, or square roots.

\subsection{Optional finite diagnostics}
The optional high-precision replay sums through $k=T^2$ at both 55 and 75 decimal digits. All 35 recorded OEIS floors agree at both precisions, with the exact omitted-tail bound \eqref{eq:numerictail}. Among $n=1,\ldots,34$, the smallest observed floor margin is approximately $0.00169001377463$, at $n=17$. This is strong consistency evidence, but the floating summation and special-function evaluations are \emph{not interval-certified}; the exact tail bound alone does not certify the finite computed sum or its floor.

A second optional route computes the completed Gaussian-line integrals directly, independently of the Abel--Plana kernel. Modes $1,2,3$ were examined for
\[
 \log x=20,40,80,120,200,400,800.
\]
The following representative first-mode values compare with the three-term approximation from \eqref{eq:mode}. The normalized error is
$(2\Re J_1-\text{three-term approximation})/A_1$.
\begin{center}
\small
\begin{tabular}{rrr}
\toprule
$\log x$ & $2\Re J_1(x)$ & Normalized error\\
\midrule
40 & $-7.25984597595149\times10^1$ & $-4.47274633407739\times10^{-5}$\\
80 & $4.33539369287231\times10^4$ & $6.72518709128962\times10^{-6}$\\
200 & $1.76009902284839\times10^{11}$ & $8.96411460870565\times10^{-8}$\\
800 & $1.04714585815476\times10^{43}$ & $-3.35563409581396\times10^{-9}$\\
\bottomrule
\end{tabular}
\end{center}
The errors are divided by the positive envelope, never by an oscillatory value. Their signs and magnitudes need not be monotone. Changing the Gaussian-line cutoff from 14 to 16 agreed to working precision; this is a stability diagnostic, not a rigorous tail or rounding certificate.

A separate symbolic implementation using SymPy agrees exactly through $d_{m,4}$ for modes $m=1,2$ with the standard-library generator. SymPy and mpmath are optional dependencies, kept separate from the mandatory exact checks.

\subsection{Archive design and limits}
The source archive contains the manuscript, original finite algebra and diagnostics, a small attributed OEIS integer fixture, reproduction instructions, and an offline deterministic builder. It excludes downloaded third-party full texts and private working reports. The builder runs the mandatory verification, compiles with a local \texttt{pdflatex}, and writes a normalized source archive. A manifest records source hashes; build checks reject invalid paths and unlisted source changes. Independent extracted rebuilds test the package rather than relying on an in-place build alone.

Reproducing the PDF requires the documented local TeX packages; no network access is needed. Mandatory algebra and tail checks require only a supported Python interpreter. Optional floating and symbolic diagnostics require their explicitly documented dependencies. Neither a successful build nor agreement of floating diagnostics supplies effective constants in the asymptotic big-$O$ terms.

\section{Source scope and further questions}\label{sec:scope}
The OEIS record attributes the sequence to Paul D. Hanna in September 2019 and the equivalent $a(n)\sim2ne^n$ to Kot\v e\v sovec two days later \cite{oeis}. The present argument proves that equivalent but concentrates on the unbounded oscillatory error concealed by it.

The classical continuous model is explicitly treated by Garrappa and Mainardi \cite{gm}; their paper includes the Volterra functions, Ramanujan integrals and prior all-orders continuous asymptotics. Abel--Plana and Stirling are standard sources \cite{dlmfap,dlmfgamma}. The general use of analytic reversion, exact-discrepancy corrections and ceiling enclosures is established methodology. Related transseries and inversion work is collected in the ProveIt exposition \cite{inversion}. No new general summation, Stirling, stationary-phase or inverse theorem is claimed here.

Bounded overlap checks examined the actual models, not just titles. A related action-accumulation model uses $\sum_{j\ge1}j^{-2}e^{-x/j}$ and Poisson--Bessel oscillations \cite{action}; another concerns late coefficients associated with A336293 \cite{late}. Neither is the present square-root gamma sampler. Other inspected comparisons involved Fabius--Rvachev operators and real Binet/Fibonacci inversion, also different models. These comparisons support a scoped description of the work, not an exhaustive literature or worldwide novelty determination.

Several questions remain beyond the proved statements:
\begin{enumerate}
\item Can the exact finite-harmonic remainder be sharpened at its boundary exponent without the arbitrary $\eps$? The present contour proof needs a strictly positive vertical decay margin.
\item Can one obtain estimates uniform in a growing number of harmonics, with a controlled transition when the mode saddle approaches the origin? Fixed-$m$ estimates alone do not answer this.
\item Can rigorous interval quadrature turn the exact contour formulas into certified numerical discrepancy and threshold enclosures, with explicit constants and an effective onset?
\item What changes for the general sampler $\sum_{k\ge0}x^{k^\rho}/\Gamma(1+k^\rho)$? The square-root case is special because the harmonic phase becomes quadratic after the natural substitution.
\end{enumerate}
These are proposed extensions, not consequences claimed from the present fixed-parameter analysis.

\begin{thebibliography}{9}
\small\raggedright
\bibitem{oeis}
The OEIS Foundation, \emph{A326805}, sequence contributed by Paul D. Hanna, 14 September 2019; asymptotic conjecture by V\'aclav Kot\v e\v sovec, 16 September 2019. Record consulted 3 October 2026. \url{https://oeis.org/A326805}.

\bibitem{gm}
R. Garrappa and F. Mainardi, \emph{On Volterra functions and Ramanujan integrals}, Analysis \textbf{36} (2016), 89--105. DOI: \href{https://doi.org/10.1515/anly-2015-5009}{10.1515/anly-2015-5009}. Author manuscript: \url{https://arxiv.org/abs/1610.01491}. In particular, Definition 1.1, Theorem 3.1 and Section 4.2.

\bibitem{dlmfap}
F. W. J. Olver et al., eds., \emph{NIST Digital Library of Mathematical Functions}, Section 2.10, Sums and Sequences, including Abel--Plana summation. \url{https://dlmf.nist.gov/2.10}.

\bibitem{dlmfgamma}
F. W. J. Olver et al., eds., \emph{NIST Digital Library of Mathematical Functions}, Section 5.11, Gamma-function asymptotic expansions and sector restrictions. \url{https://dlmf.nist.gov/5.11}.

\bibitem{inversion}
V. Reshetnikov, \emph{Transseries and Inversion}, ProveIt repository exposition, consulted 3 October 2026. \href{https://github.com/VladimirReshetnikov/ProveIt/tree/main/Analysis/Transseries/docs/series-and-transseries/Transseries_And_Inversion}{Online exposition}.

\bibitem{action}
V. Reshetnikov, \emph{Action Accumulation Nonlinear Inversion}, ProveIt repository exposition, consulted 3 October 2026. \href{https://github.com/VladimirReshetnikov/ProveIt/tree/main/Analysis/Transseries/docs/series-and-transseries/Action_Accumulation_Nonlinear_Inversion}{Online exposition}.

\bibitem{late}
V. Reshetnikov, \emph{Late Growth Bessel Counting Coefficients}, ProveIt repository exposition, consulted 3 October 2026. \href{https://github.com/VladimirReshetnikov/ProveIt/tree/main/Analysis/Transseries/docs/series-and-transseries/Late_Growth_Bessel_Counting_Coefficients}{Online exposition}.
\end{thebibliography}
\end{document}
